% Part: propositional-logic
% Chapter: propositional-logic
% Section: soundness

\documentclass[../../../include/open-logic-section]{subfiles}

\begin{document}

\olfileid{pl}{prp}{snd}

\olsection{विधानात्मक तर्कशास्त्राची निर्दोषता}

\begin{thm}[निर्दोषता]
\ollabel{thm:soundness}
$\Gamma\Proves!A$ असेल, तर
$\Gamma \Entails !A$.
\end{thm}

\begin{proof}
प्रमेयांवर आगमन करू. $!A$ स्वयंसिद्धक असेल, तर $\Entails
!A$; म्हणून विशेषतः $\Gamma \Entails!A$. तसेच $!A \in \Gamma$
असेल, तर $\Gamma \Entails!A$. आगमनात्मक पायरीसाठी $!C$
आणि $!C\lif!A$ यांपासून \emph{विध्यात्मक अनुमानाने} $!A$
मिळाले असे माना. तसेच आगमनाच्या गृहीतकानुसार प्रमेय
$!C\lif !A$ आणि $!C$ यांसाठी खरे आहे असे माना.
\olref[sem]{prop:semanticalfacts} मधील तिसऱ्या मुद्द्यावरून
अपेक्षित निष्कर्ष मिळतो.
\end{proof}

\begin{cor}
$\Gamma$ पूर्ततायोग्य असेल, तर $\Gamma$ सुसंगत असतो. म्हणून
विधानात्मक तर्कशास्त्र सुसंगत आहे.
\end{cor}

\end{document}
